\documentclass[11pt]{article}
\usepackage[T1]{fontenc}
\usepackage{amsmath}
\usepackage{amssymb}
\usepackage{amsthm}
\usepackage[hidelinks]{hyperref}
\newtheorem{theorem}{Theorem}[section]
\newtheorem{lemma}[theorem]{Lemma}
\newtheorem{corollary}[theorem]{Corollary}
\theoremstyle{definition}
\newtheorem{definition}[theorem]{Definition}
\begin{document}
\begin{center}
{\LARGE Propositional logic}
\end{center}
\section{Conjunction}
\label{ll:main:conjunction}
\begin{theorem}
\label{ll:main:and-comm}
For every proposition \(p\) and proposition \(q\), if \(p\) and \(q\), then \(q\) and \(p\).
\end{theorem}
\begin{proof}
Assume \(h\).
Consider the cases of \(h\).
Case \(\operatorname{andintro}\) with \(hp\), \(hq\):
\begin{quote}
Branch \(1\):
\begin{quote}
The goal follows from \(hq\).
\end{quote}
Branch \(2\):
\begin{quote}
The goal follows from \(hp\).
\end{quote}
\end{quote}
\end{proof}
\begin{theorem}
\label{ll:main:and-assoc}
For every proposition \(p\) and proposition \(q\) and proposition \(r\), \(p\) and \(q\) and \(r\) if and only if \(p\) and \(q \land r\).
\end{theorem}
\begin{proof}
Branch \(1\):
\begin{quote}
Assume \(h\).
Consider the cases of \(h\).
Case \(\operatorname{andintro}\) with \(hpq\), \(hr\):
\begin{quote}
Consider the cases of \(hpq\).
Case \(\operatorname{andintro}\) with \(hp\), \(hq\):
\begin{quote}
Branch \(1\):
\begin{quote}
The goal follows from \(hp\).
\end{quote}
Branch \(2\):
\begin{quote}
Branch \(1\):
\begin{quote}
The goal follows from \(hq\).
\end{quote}
Branch \(2\):
\begin{quote}
The goal follows from \(hr\).
\end{quote}
\end{quote}
\end{quote}
\end{quote}
\end{quote}
Branch \(2\):
\begin{quote}
Assume \(h\).
Consider the cases of \(h\).
Case \(\operatorname{andintro}\) with \(hp\), \(hqr\):
\begin{quote}
Consider the cases of \(hqr\).
Case \(\operatorname{andintro}\) with \(hq\), \(hr\):
\begin{quote}
Branch \(1\):
\begin{quote}
Branch \(1\):
\begin{quote}
The goal follows from \(hp\).
\end{quote}
Branch \(2\):
\begin{quote}
The goal follows from \(hq\).
\end{quote}
\end{quote}
Branch \(2\):
\begin{quote}
The goal follows from \(hr\).
\end{quote}
\end{quote}
\end{quote}
\end{quote}
\end{proof}
\section{Disjunction}
\label{ll:main:disjunction}
\begin{theorem}
\label{ll:main:or-comm}
For every proposition \(p\) and proposition \(q\), if \(p\) or \(q\), then \(q\) or \(p\).
\end{theorem}
\begin{proof}
Assume \(h\).
Consider the cases of \(h\).
Case \(\operatorname{orinl}\) with \(hp\):
\begin{quote}
Select the right alternative.
The goal follows from \(hp\).
\end{quote}
Case \(\operatorname{orinr}\) with \(hq\):
\begin{quote}
Select the left alternative.
The goal follows from \(hq\).
\end{quote}
\end{proof}
\begin{theorem}
\label{ll:main:or-elim}
For every proposition \(p\) and proposition \(q\) and proposition \(r\), if \(p\) or \(q\), then if if \(p\), then \(r\), then if if \(q\), then \(r\), then \(r\).
\end{theorem}
\begin{proof}
Assume \(hpq\), \(hpr\), \(hqr\).
Consider the cases of \(hpq\).
Case \(\operatorname{orinl}\) with \(hp\):
\begin{quote}
Apply \(hpr\).
The goal follows from \(hp\).
\end{quote}
Case \(\operatorname{orinr}\) with \(hq\):
\begin{quote}
Apply \(hqr\).
The goal follows from \(hq\).
\end{quote}
\end{proof}
\section{Negation}
\label{ll:main:negation}
\begin{theorem}
\label{ll:main:double-negation-intro}
For every proposition \(p\), if \(p\), then not not \(p\).
\end{theorem}
\begin{proof}
Assume \(hp\).
Assume \(hn\).
Apply \(hn\).
The goal follows from \(hp\).
\end{proof}
\begin{theorem}
\label{ll:main:de-morgan}
For every proposition \(p\) and proposition \(q\), if not \(p \lor q\), then not \(p\) and not \(q\).
\end{theorem}
\begin{proof}
Assume \(h\).
Branch \(1\):
\begin{quote}
Assume \(hp\).
Apply \(h\).
Select the left alternative.
The goal follows from \(hp\).
\end{quote}
Branch \(2\):
\begin{quote}
Assume \(hq\).
Apply \(h\).
Select the right alternative.
The goal follows from \(hq\).
\end{quote}
\end{proof}
\begin{theorem}
\label{ll:main:explosion}
For every proposition \(p\) and proposition \(q\), if \(p\) and not \(p\), then \(q\).
\end{theorem}
\begin{proof}
Assume \(h\).
Consider the cases of \(h\).
Case \(\operatorname{andintro}\) with \(hp\), \(hn\):
\begin{quote}
The goal follows from \(\operatorname{absurd}(hp, hn)\).
\end{quote}
\end{proof}
\section{Equivalence}
\label{ll:main:equivalence}
\begin{theorem}
\label{ll:main:iff-refl}
For every proposition \(p\), \(p\) if and only if \(p\).
\end{theorem}
\begin{proof}
Branch \(1\):
\begin{quote}
Assume \(h\).
The goal follows from \(h\).
\end{quote}
Branch \(2\):
\begin{quote}
Assume \(h\).
The goal follows from \(h\).
\end{quote}
\end{proof}
\begin{theorem}
\label{ll:main:iff-symm}
For every proposition \(p\) and proposition \(q\), if \(p\) if and only if \(q\), then \(q\) if and only if \(p\).
\end{theorem}
\begin{proof}
Assume \(h\).
Consider the cases of \(h\).
Case \(\operatorname{iffintro}\) with \(mp\), \(mpr\):
\begin{quote}
Branch \(1\):
\begin{quote}
The goal follows from \(mpr\).
\end{quote}
Branch \(2\):
\begin{quote}
The goal follows from \(mp\).
\end{quote}
\end{quote}
\end{proof}
\section{Classical reasoning}
\label{ll:main:classical}
\begin{theorem}
\label{ll:main:double-negation-elim}
For every proposition \(p\), if not not \(p\), then \(p\).
\end{theorem}
\begin{proof}
Assume \(hnn\).
Consider the cases of \(\operatorname{em}(p)\).
Case \(\operatorname{orinl}\) with \(hp\):
\begin{quote}
The goal follows from \(hp\).
\end{quote}
Case \(\operatorname{orinr}\) with \(hn\):
\begin{quote}
The goal follows from \(\operatorname{absurd}(hn, hnn)\).
\end{quote}
\end{proof}
\end{document}
